% status: FULL
\chapter{CPU and SIMD}\label{app:cpu}
\begin{ChapterIntent}
A CPU is both a numerical target and the controller of an accelerator
service. We will follow one array through a scalar loop, vector lanes and
several threads. The matrix products of Chapter~\ref{app:matmul} are our
starting point; instruction-specific performance belongs to the later
systems chapters' measurement records.
\end{ChapterIntent}

\section{Elements are addressed; lines are transferred}
An array access names an element, but memory hardware commonly transfers
a cache line. Sequential access can use most of it. A large stride may
transfer many bytes for each useful element. On an illustrative machine
with 64-byte lines and four-byte elements, a line holds sixteen elements.
Reading one element per distinct line uses one sixteenth of its payload
before reuse. This is not the actual M1 line size.

Row-major traversal works well because neighbouring column indices name
neighbouring elements. Transposing a view changes strides without
necessarily moving data. A loop that ignores those strides can be wrong;
one that respects them can still be slow. Chapter~\ref{app:tensors} separates these
failures. Neither a contiguous allocation nor a matrix-shaped API proves
that a chosen traversal has locality.

\section{Vector lanes and the tail}
Single instruction, multiple data (SIMD) applies an operation across
several lanes. It helps when data can feed those lanes and the compiler
can establish safe accesses. Branching, aliasing, strides and reductions
can complicate that proof. An intrinsic makes an operation explicit but
does not make the surrounding memory access efficient.

For an illustrative eight-lane vector processing nineteen elements, two
full vectors leave three elements. A scalar cleanup or correctly masked
operation handles that tail. Loading eight elements from a three-element
region without a valid memory contract is not an optimization. Test
lengths just below and above vector boundaries, and test misaligned views
if the interface allows them.

NEON, AVX and matrix extensions differ in width, operations and dispatch
requirements. Check the actual target and runtime support rather than
inferring support from the operating system. A portable scalar reference
is valuable even when the optimized path is architecture-specific.

\section{Reduction order is part of the result}
Several independent accumulators can expose instruction-level parallelism.
Combining them changes summation order. Real-number algebra allows this;
floating-point arithmetic need not produce identical bits. A random seed
does not repair an altered arithmetic result. Chapter~\ref{app:numerics}
explains the numerical contract.

First compare the output against an independent reference. Then inspect
vectorization reports or generated instructions. Finally measure the
intended data size and thread count. A tiny array in cache answers a
different question from a checkpoint larger than the last-level cache.

\section{Threads share resources}
Threads may share caches and a memory interface. Once aggregate traffic
saturates the interface, more threads can add overhead without increasing
throughput. False sharing occurs when threads write different values in
the same coherence line: logically independent updates still trigger
ownership transfers. Separate per-thread accumulators and combine them
deliberately.

Record affinity, the cores actually used, and nested library thread counts.
Giving every nested layer all available cores can oversubscribe the
machine. On a NUMA system, page placement and the accessing core determine
which memory path is used. Establish locality before calling one result
the machine's memory bandwidth.

\section{The CPU also runs the service}
A CPU path can win for tiny operations whose accelerator submission costs
more than their execution. The same CPU may tokenize, enforce grammars,
schedule batches and service streams. A faster isolated matrix kernel can
therefore worsen a server if it occupies control-plane cores.

Predict the shared resource, measure a thread-count sweep, and retain a
one-thread baseline. Never multiply a vector-width ratio by a core-count
ratio and present the product as measured speedup. The roofline and
launch-overhead chapters demonstrate how to test that hypothesis.

\section{Walk through a nineteen-element loop}
Take the integers one through nineteen, stored here as numerical values
whose exact sum is easy to check. A scalar loop starts at zero and adds
one value on each iteration. Its answer is
$19\cdot20/2=190$. This is our oracle before considering vectorization.

An illustrative eight-lane implementation reads values one through eight,
then nine through sixteen. Each lane accumulates two values. Reducing
the eight lane sums produces 136, the sum of one through sixteen.
The remaining three values are 17, 18 and 19, which contribute 54.
Adding the tail restores 190.

\EngineFigure{foundation-simd-tail}{Nineteen elements processed with an
illustrative eight-lane vector width. Two complete groups are followed by
a three-element tail. The grouping changes how work is issued, not which
elements belong to the sum.}{fig:foundation-simd-tail}

The arithmetic exposes a common defect. A loop bounded by the number of
complete vectors can produce 136 without any crash or invalid number.
An input size divisible by eight would hide it. Test lengths seven,
eight, nine, fifteen, sixteen, seventeen and nineteen, as well as zero
when the interface permits an empty input.

Run the small arithmetic lesson:
\begin{verbatim}
python3 code/foundations/lesson.py cpu
\end{verbatim}
It reports two complete groups, three tail elements and sum 190.
It does not claim that Python executed SIMD instructions. The lesson
checks the indexing contract; the Rust matrix kernels and their
benchmarks are where we inspect actual execution.

\section{A dependency can keep lanes waiting}
In a scalar reduction, each new sum depends on the previous sum.
Even if the processor can issue several arithmetic instructions,
that dependency chain can limit progress. Independent accumulators
allow work to proceed before earlier additions complete.

For example, accumulate alternating elements into two partial sums,
then add the partial sums at the end. Integer examples can make this
look identical to the scalar loop. With floating-point values, the
different grouping can change rounding. The optimization therefore
needs both a performance test and a numerical contract.

This distinction helps when reading compiler output. Finding a vector
instruction confirms that some work was vectorized. It does not show
that the complete loop is fast, that memory supplies enough values,
or that its numerical result meets the application contract.
Inspecting generated code and timing the intended workload answer
different questions.

\section{Data reuse is often more valuable than wider lanes}
Suppose several input rows are multiplied by the same weight row.
A loop that finishes one input at a time may reload that weight row
for every input. A loop that keeps the weight row available while
updating several outputs can reuse each fetched value.
The scalar products are the same; the data movement changes.

This is the bridge from a vector operation to a small matrix tile.
The tile size must fit useful values and accumulators into available
registers or caches. Making it larger without limit eventually causes
spills, extra loads or capacity misses. A tile is a trade between
reuse and local storage, not a synonym for maximum size.

A useful experiment varies the number of input rows while keeping the
weight matrix fixed. First predict whether reuse should improve.
Then measure under controlled cache and thread conditions. If the
curve changes abruptly, inspect dispatch thresholds or blocking choices
before concluding that the hardware itself has a discontinuity.
The later matrix-multiplication chapter performs that investigation.

\section{Threads need independent outputs}
Parallelizing matrix multiplication over output rows gives workers
different destinations. Each worker can read the shared input and
weights while writing its own output region. The program still needs
to ensure those regions do not overlap.

Parallelizing one reduction is harder. Workers produce partial sums
that must be combined, introducing synchronization and a new reduction
order. For a very small reduction the coordination can cost more
than the work it distributes. Begin with the simpler ownership split
and measure whether the problem is large enough to pay for threads.

Even disjoint scalar outputs can share a cache line. Repeated writes
to that line by different cores can trigger coherence traffic.
Padding or rearranging per-thread state can help, but padding increases
memory footprint. Do not pad everything by habit. Identify the shared
line, show the update pattern and measure the changed layout.

Thread count is not the same as useful core count. A program can
launch more runnable threads than available execution resources.
Nested parallel libraries can amplify this problem: several outer
tasks each launch their own full team of inner workers. Record both
levels when interpreting a speedup or slowdown.

\section{Read a CPU benchmark in layers}
Check correctness first. Then name the timed data: an in-cache array,
a streaming matrix, or a request including allocation and tokenization.
Record repeated input reuse, excluded costs, thread count and competing
work. A valid warm-cache measurement does not directly estimate
checkpoint streaming; the later roofline chapter measures separate
boundaries for that reason.

Include the controller's work when evaluating a server. Occupying every
core can delay token delivery, admission or cancellation while improving
an isolated kernel. The measurements have different scopes.

\section{A practical first optimization}
Start from a checked scalar loop with explicit indexing, ownership
and tails. Reuse data before adding target-specific instructions.
Measure one thread before many; retain earlier results and rerun
correctness after each change. For larger tiles the same questions
remain: what is reused, where does it live, which operations are
independent, and which numerical contract must survive?
\section{Turn a stride into a set of transferred lines}
Let $\mathbf{A}\in\mathbb{R}^{16\times16}$ be a row-major F32
matrix, aligned to a line boundary on our illustrative 64-byte-line
machine. There are 256 values and 1,024 payload bytes. Element
$(i,j)$ begins at the byte address
\[
 a(i,j)=a_0+4(16i+j),
\]
where $a_0$ is the allocation's base address. The line index is
$\lfloor a(i,j)/64\rfloor$. We count distinct line indices touched
by an access pattern, not source-code loads.

A complete row uses sixteen adjacent words in one line. The
whole matrix spans sixteen lines. Now read only column zero.
Successive values lie at offsets $0,64,128,\ldots,960$.
Those sixteen F32 values contain 64 useful bytes, but name the
same sixteen distinct lines. Under the stated empty-cache,
single-read counting model, the line payload is 1,024 bytes
and useful-byte fraction is $64/1024=1/16$.

\EngineFigure{foundation-cpu-lines}{Two access patterns over the same
illustrative aligned matrix. Dark first-column cells identify the
strided pattern; a row consumes a whole line. The count names distinct
64-byte lines, not measured misses, DRAM bytes or the actual line
size of the machine running the lab.}{fig:foundation-cpu-lines}

Two details prevent this drawing from becoming a false performance
prediction. First, lines might already reside in a cache, and
another pass can reuse them. Second, fetching one line can overlap
with other outstanding requests. Latency and bandwidth constrain
different parts of that process. The address set explains why
one traversal can waste payload; it cannot prove a sixteenfold
slowdown.

The executable \texttt{matrix\_lines} lesson computes these sets
from byte addresses. It also tests a four-byte access beginning
at byte 62: the access spans two 64-byte lines. Looking only at
the starting address would undercount it. This is a useful small
analogy for vector loads crossing a boundary, though a legal
hardware load and an allocation boundary are separate contracts.

\section{A cache is neither a tensor nor a permanent allocation}
The tensor API describes numerical coordinates and ownership.
A cache holds recently used lines according to a hardware
policy. Keeping an array allocated does not promise that its
lines remain in a particular cache. A sequence of source-level
loads does not establish how often DRAM was read.

To reason about reuse, name the working set: the values needed
before they are used again. A matrix tile may include part
of a weight row, several input rows and partial output sums.
Its byte footprint is the sum of those live objects, not only
the weights. A tile that exceeds a useful storage level can
evict data before reuse or force register values to spill.
This is why reducing one load in the source can add several
others after compilation.

For a concrete example, eight output accumulators held as F32
use 32 bytes of numerical state. This number is not a register
allocation guarantee: address variables, loop state and compiler
temporaries also need resources, and the target has its own
allocation rules. Raising a tile's output count increases its
reuse opportunity and its live state together. Look at both
before choosing a block size.

Alignment also has several meanings. A pointer can be aligned
for an F32 value but not for a sixteen-byte vector or a cache
line. An allocation can be well aligned while a sliced view
begins four bytes later. Our native lab accepts float-aligned
pointers and tests offset zero and offset one; it does not
declare a stronger alignment promise to the compiler. A
production interface should say whether it supports such views.

\section{Read a real four-lane loop}
The earlier eight-lane drawing taught the tail contract without
selecting an instruction set. We now execute a different, concrete
path: four F32 lanes using AArch64 NEON. Each vector represents
four independent products. Arm's intrinsic reference describes
the load, multiply, add and horizontal reduction used here
\cite{arm2026neon}. The scalar path remains available on other
targets and reports its identity.

\begin{minipage}{\linewidth}
\CodeLines{../code/foundations/cpu-dot.c}{14}{26}
\end{minipage}

The vector loop condition is \texttt{n - i >= 4}. At entry,
$0\leq i\leq n$, and each iteration advances by four only
when at least four values remain. That invariant makes the
unsigned subtraction safe and prevents an out-of-range vector
load. After the loop, $n-i$ is zero, one, two or three.
The scalar cleanup consumes precisely those remaining values.
The condition is not an invitation to load extra elements
and discard them later: their storage still has to be valid.

At $n=19$, the native path consumes four full vectors, covering
indices zero through fifteen, then three scalar values. Four
lanes is not the same grouping as the earlier eight-lane example,
but both cover the same nineteen indices. An implementation
can change grouping without changing the input contract.

Each lane accumulates positions separated by four. Lane zero
holds products at indices $0,4,8,\ldots$; lane one holds
$1,5,9,\ldots$. The horizontal reduction combines lane sums,
and the tail is added afterwards. This is not the left-to-right
scalar order. For the bounded dyadic fixture below, both are
exact. For general weights, Chapter~\ref{app:numerics} supplies
the reason to compare under a declared error contract.

\EngineFigure{foundation-neon-reduction}{Eight products through the
four-lane teaching loop. Vertical arrows show contributions to one
lane's subsequence; the final operation combines four lane sums.
The diagram omits loads and depicts grouping, not instruction
cycles or measured speed.}{fig:foundation-neon-reduction}

\section{Choose inputs whose oracle is independent}
The native lab does not make the scalar float loop its only
oracle. For index $i$, it chooses
\[
 a_i=\frac{(i\bmod13)-6}{4},\qquad
 b_i=\frac{(i\bmod7)-3}{2}.
\]
The residues are interpreted as signed integers before subtraction.
Thus each product is an integer divided by eight. An independent
integer loop sums the numerator, then divides once by eight
to form the expected answer.

With at most 257 values, each numerator's magnitude is at
most $6\cdot3=18$, so even the sum of absolute numerator
magnitudes is at most $257\cdot18=4626$. These eighths are
exactly representable throughout the scalar and vector
accumulations in F32. We may therefore require exact equality
for this fixture without pretending that all floating-point
dot products are exact.

The lab checks every length from zero to 257 at offsets zero
and one: $258\cdot2=516$ cases. The allocation contains 258
values, so the final offset-one case ends exactly at the
allocation boundary. Short lengths exercise no vector loop;
lengths surrounding every multiple of four exercise different
tails. The offset-one path tests a valid F32 slice that need
not have the vector alignment of the base pointer.

Remove the scalar tail in a scratch copy. Some divisible-by-four
cases still pass, but non-divisible lengths need not. The
record includes a deliberately tail-free variant whose
length-one failure is detected. Predict that failure from
the first product before running it. A test distribution
containing only aligned multiples of four would have missed
the bug.

\section{Separate a compiled instruction from a benchmark}
Build and run from the repository root:
\begin{verbatim}
mkdir -p build/foundations-machine
cc -std=c11 -O2 -ffp-contract=off -Wall -Wextra -Werror \
  code/foundations/cpu-dot.c -o build/foundations-machine/cpu-dot
build/foundations-machine/cpu-dot
\end{verbatim}
The result names either the NEON path or the scalar fallback
and reports 516 successful cases. The local run used the NEON
path; its environment and result are recorded separately.
No elapsed time is reported.

Generate assembly with the same target and numerical flags:
\begin{verbatim}
cc -std=c11 -O2 -ffp-contract=off -S \
  code/foundations/cpu-dot.c -o build/foundations-machine/cpu-dot.s
\end{verbatim}
Find the function \texttt{lane\_dot}, not only an instruction
somewhere in the file. The inspected AArch64 output contains
four-lane \texttt{fmul} and \texttt{fadd}, followed by
\texttt{faddp} operations for reduction. A compiler can select
different load mnemonics from those suggested by an intrinsic's
documentation. Read the operand widths and function context,
not a string match alone.

The contraction flag keeps a multiply followed by add from
being silently treated as our fused path. It is not a blanket
determinism guarantee across compilers or processors. This lab
avoids reassociation-enabling fast-math flags. If a later kernel
uses FMA deliberately, its reference should describe that
numerical choice instead of comparing to the wrong graph.

Assembly inspection establishes that vector instructions were
generated for this build. It does not establish how many
instructions retired per second, whether loads hit cache or
whether the server improves. Those require an experiment with
an explicit workload and timing scope. Keep the generated
code and runtime evidence as separate layers.

\section{Partition writes before adding worker threads}
For $\mathbf{Y}=\mathbf{X}\mathbf{W}^{\mathsf T}$, suppose
$\mathbf{X}\in\mathbb{R}^{B\times d}$ and
$\mathbf{W}\in\mathbb{R}^{m\times d}$. An output element
$(r,j)$ is a reduction over $d$. Splitting output rows
or ranges of output columns gives each worker an exclusive
destination while inputs remain read-only. There is no
mathematical need to combine workers' partial results if
each complete reduction has one owner.

Splitting the reduction axis instead produces several partial
values for the same destination. These must be merged in
some defined way. Concurrent unsynchronized updates are a
data race, not an alternative reduction policy. Locking
each update establishes mutual exclusion but can serialize
the useful work and still vary the order of additions.
Separate per-worker partial buffers and a deliberate merge
make the ownership and numerical order easier to inspect.

Exclusive destinations can nevertheless share a coherence
line. Imagine four-byte counters at offsets zero and four
in one illustrative 64-byte line. Two workers repeatedly
update different counters. They do not conflict at the
language-level object boundary, but a coherent machine may
transfer ownership of the same line between cores. Place
the counters in separate lines and the ownership interference
can change. It is a hypothesis to measure, not a promise
that every padded structure is faster.

A NUMA topology adds placement. A worker can own an output
range while its pages reside behind another memory controller.
Ownership correctness does not imply locality. Likewise,
four outer tasks each starting eight library workers create
32 potentially runnable numerical workers before accounting
for control work. They do not create 32 physical cores.
State outer concurrency and inner thread counts together.

\section{Reuse weights across a small batch}
For one input row and one weight row of length $d$, a dot
product performs $d$ multiplications and $d$ additions under
our two-FLOPs-per-multiply-add convention. Ignoring the small
output, F32 input and weight traffic would be $8d$ bytes
if both rows were fetched anew at this boundary. Arithmetic
intensity is then approximately $2d/(8d)=0.25$ FLOP/byte.

Now update four input rows while reusing the weight row.
The arithmetic is $8d$ FLOPs and the named-boundary input
traffic is $(4+1)4d=20d$ bytes, giving 0.4 FLOP/byte.
The reusable object is the weight row, not the four input
rows magically disappearing. If input rows are already
resident at the boundary, the denominator changes again.
Always name which storage boundary the traffic count describes.

A small tile exposes that reuse by keeping several output
accumulators live while stepping through the weight row.
The larger matrix case extends the same idea in more than
one dimension. There are costs: more accumulators, more
address state, boundary handling and possibly a different
reduction order. Wider SIMD and tiling solve related but
different problems. SIMD processes several elements per
instruction; tiling tries to reuse fetched values.

\section{The controller needs a budget too}
When numerical workers share a CPU with an accelerator service,
ask what happens to cancellation, tokenization and output
delivery while the kernel runs. A delayed cancellation can
retain KV state and spend more accelerator steps even if the
CPU kernel became faster in isolation. This is a resource
interaction, not a reason to forbid CPU inference.

A first experiment should compare one numerical worker,
several workers and the normal service configuration. Keep
the same input and error contract. Record whether token
delivery and admission are inside the timing boundary.
Then identify whether the limit is compute, memory supply,
coordination or control-plane contention. The later request
lifecycle and roofline chapters make these measurements.
This chapter has established the addresses, lanes and writers
needed to interpret them without manufacturing a benchmark.

\begin{ConstraintLedger}
Memory bandwidth & buys & Contiguous access and weight reuse increase useful work per fetched line.\\
Compute & buys & Independent lanes and accumulators expose parallel operations.\\
Memory capacity & spends & Live tiles and per-worker state consume registers and cache footprint.\\
Correctness & risks & Tails, aliasing, races and changed reduction order need separate checks.\\
Latency & risks & Numerical workers can delay the service's control work.\\
\end{ConstraintLedger}
\section{Worked problems}
\subsection*{1. The missing tail}
\textbf{Problem.} An eight-lane loop sums only the first sixteen of the
values one through nineteen. How wrong is its answer?
\textbf{Solution.} It returns 136 instead of 190, omitting $17+18+19=54$.
\subsection*{2. Useful bytes}
\textbf{Problem.} With illustrative 64-byte cache lines and F32 values,
what payload fraction is used by one value per distinct line before reuse?
\textbf{Solution.} $4/64=1/16$, or 6.25 percent. This is useful payload
per transferred line, not a universal cache-miss or speedup prediction.
\subsection*{3. Choose the split}
\textbf{Problem.} Why begin by splitting a matrix product across output
rows rather than splitting every dot-product reduction?
\textbf{Solution.} Distinct output rows admit independent writers without
a partial-sum merge. Splitting a reduction adds coordination and changes
summation order; it can help at large sizes but needs separate evidence.

